\section{Qualidade de Software e Testes}
\label{sec:qualidade}

A garantia de qualidade é um pilar fundamental na engenharia de software, especialmente em sistemas distribuídos e assíncronos como o proposto neste trabalho. A complexidade inerente a uma arquitetura orientada a eventos, com múltiplos componentes desacoplados, exige uma abordagem disciplinada e multifacetada para assegurar a corretude, a robustez e a confiabilidade do sistema.

Neste contexto, a estratégia de qualidade adotada não se limitou apenas à verificação funcional do software, mas foi integrada ao próprio processo de design e desenvolvimento. Foram empregadas práticas de análise estática de código, como 	extit{linting} e verificação de tipos, para garantir a consistência e a legibilidade. Além disso, uma estratégia de testes automatizados em múltiplas camadas foi implementada para validar o sistema em diferentes níveis de abstração, desde a lógica de negócio de componentes individuais até o fluxo de processamento completo da aplicação. A subseção a seguir detalha o plano de testes executado, que constitui a base para a validação da qualidade e da corretude da plataforma desenvolvida.
